\chapter{Feature Completion as Positive and Negative Proof}
\label{ch:feature-discipline}

A request to show a message, list source, accept a transaction, render a diagram
or produce an MP4 states a desirable observation. The change can also introduce
stale state, disclosure, resource exhaustion, weaker authentication, corruption
or success without an artifact. Completion therefore requires a paired argument:

\begin{enumerate}
  \item a \emph{positive invariant} states what the system must do; and
  \item applicable \emph{negative invariants} state what the system must never do.
\end{enumerate}

Review before declaring completion, in proportion to the boundary crossed:
button spacing rarely needs a cryptographic audit; filesystem access needs
authorization and path analysis. A material risk without an accepted control
blocks completion.

\section{The logical shape of a feature}

Let $good_f$ mean the observable contract of feature $f$, and let
$bad_{f,1},\ldots,bad_{f,n}$ name reachable failures introduced or exposed by
that feature.  A common safety-shaped completion condition is

\AccessibleEquation{On every reachable future the feature contract holds; on every reachable future none of the identified bad outcomes holds.}{\[
  AG\,good_f
  \;\land\;
  \bigwedge_{i=1}^{n} AG\,\neg bad_{f,i}.
\]}

The first conjunct is positive: the requested behavior remains true throughout
every permitted execution.  Each remaining conjunct is negative: no permitted
execution reaches a forbidden observation.  The equivalent counterexample form

\AccessibleEquation{A forbidden outcome is absent exactly when there does not exist a future that eventually reaches it.}{\[
  AG\,\neg bad \;\equiv\; \neg EF\,bad
\]}

is operationally useful.  A negative test tries to construct the $EF\,bad$
witness: another user's project path, a stale response after a project switch,
an oversized upload, a duplicate response, or a secret in a log.  Finding such
a witness refutes the claimed completion.

Not every feature is purely invariant.  ``A submitted request eventually gets a
terminal response'' is a liveness obligation such as
$AG(submitted\rightarrow AF(complete\lor failed))$.  A feature review must not
replace liveness with safety or hide fairness and availability assumptions.
Positive and negative requirements classify intent; the temporal operator must
still match the claim.

\section{Stability review before security review}

Before editing code, identify the existing invariant and review the applicable
ways an environment change could break it:

\begin{description}
  \item[Regression.] Which accepted behavior, file format, route, proof, or UI
    state can this change invalidate?  Preserve conforming behavior and add a
    regression test at the narrowest durable boundary.
  \item[Failure.] What happens when a dependency, renderer, compiler, provider,
    filesystem operation, or network request fails?  Failure must become an
    explicit terminal observation rather than a false success or endless wait.
  \item[Concurrency.] Can a late response, simultaneous session, queued prompt,
    or overlapping write be attributed to the wrong project or revision?
    Correlation identity and revision checks must survive reordering.
  \item[Resources.] What bounds memory, bytes, queue occupancy, recursion,
    process time, media dimensions, and response size?  A new feature must not
    silently turn a finite resource contract into an unbounded one.
  \item[Durability and recovery.] Which state is authoritative after restart or
    partial failure?  Atomic writes, rebuildable indexes, and project-local
    ownership make recovery inspectable.
  \item[Compatibility.] Which clients, stored revisions, generated artifacts,
    and build environments must continue to work?  Version or migrate a changed
    representation rather than reinterpreting old bytes silently.
\end{description}

Retain concrete invariants at the changed boundary, briefly recording when a
dimension is non-material.

\section{Security as forbidden observable behavior}

``A user never reads a project they do not own'' states protected behavior;
``use authentication'' only names a mechanism. Before selecting a control,
identify:

\begin{enumerate}
  \item the trust boundary crossed by the feature;
  \item the asset being protected;
  \item the attacker's relevant capability; and
  \item the forbidden observable outcome.
\end{enumerate}

Applicable forbidden behaviors include unauthorized access, cross-user or
cross-project disclosure, secret persistence or logging, injection, path or
symlink escape, replay, request/response confusion, privilege expansion,
executable media, resource exhaustion and fail-open behavior. State testable
invariants, such as ``a source path absent from the authorized project inventory
never returns bytes.''

Map each material negative invariant to a control---canonical-path confinement,
authenticated project identity, an allowlisted inventory, a byte limit or an
atomic revision comparison---and a negative test or proof. A mechanism alone
does not establish the forbidden transition's absence.

\section{The feature-completion ledger}

For each feature, retain the following compact traceability record in the typed
requirements and implementation mapping:

\begin{center}
\begin{tabular}{p{0.19\linewidth}p{0.72\linewidth}}
\toprule
Field & Required content \\
\midrule
Prompt & The stakeholder request whose completion is being claimed. \\
Positive claim & Observable always or liveness property, with actor, scenario,
boundary, and revision. \\
Negative claims & Applicable never or always-not properties; omitted families
are explicitly judged non-material. \\
Controls & Concrete implementation mechanisms mapped one-to-one or
many-to-one to the negative claims. \\
Evidence & Positive conformance tests plus negative tests, proofs, generated
traces, and named checks. \\
Residual risk & Accepted limitation, assumption, deferred mitigation, or an
explicit completion blocker. \\
\bottomrule
\end{tabular}
\end{center}

Both sides need evidence: forbidding failures does not demonstrate the requested
feature, and a successful demonstration does not rule out disclosure or corruption.

\section{Worked example: the project Source view}

Consider the request ``show implementation files in the Source tab.''  The
positive claim is that selecting an authorized project lists its implementation
files and selecting one readable UTF-8 file displays its content.  Project
switching clears the old selection, so a late response cannot repopulate the new
project with stale content.

The feature crosses a filesystem boundary.  Project source, ignored account
records, repository metadata, and provider credentials are protected assets.  A
user can choose a URL path and can guess names that the visible tree never
listed.  Applicable negative invariants therefore include:

\begin{itemize}
  \item a user never reads another user's project source;
  \item a guessed ignored or metadata path never returns bytes;
  \item a path or symlink never escapes the authorized source root;
  \item a binary or oversized file never becomes an unbounded text response;
  \item a project switch never retains the previous project's source state; and
  \item an empty source directory never accompanies a claim that implementation
        succeeded.
\end{itemize}

Controls are an authenticated explicit project ID, canonical root resolution,
exact source-inventory membership, UTF-8 reads, a finite response limit and
project-keyed React state. Negative tests request another account's file and a
guessed \texttt{.git/config}; both return not-found without revealing existence.
The positive test creates, lists and opens a project-local file and checks its
exact content.

\section{Worked example: chat auto-scroll}

Even a small convenience has stability boundaries.  The positive invariant is
that initial conversation load and every newly rendered user, assistant, or
failure message reveal the bottom of the chat message pane.  The implementation
must scroll that pane rather than the entire page, must not move the independent
project sidebar, and must reset with the project-keyed chat component.  No new
authorization boundary is crossed, so a full threat model would be noise; the
project-isolation and stale-response invariants remain material.

\section{When the result is not complete}

An unobservable or untested positive claim, missing durable negative requirement,
uncontrolled risk or material control without a negative test or proof leaves
the ledger incomplete. So does a liveness claim missing its fairness or
availability assumptions, failed or skipped validation, or a check of a
substitute artifact. Report the blocker and return to the relevant gate;
stakeholder-visible residual risk requires acceptance there. A green build or
persuasive diagram does not fill these gaps.
